% 星际行星（综述）
% license CCBYSA3
% type Wiki

本文根据 CC-BY-SA 协议转载翻译自维基百科\href{https://en.wikipedia.org/wiki/Rogue_planet}{相关文章}。

流浪行星（rogue planet），亦称自由漂浮行星（free-floating planet，FFP）或孤立的行星质量天体（isolated planetary-mass object，iPMO），是指一种具有行星质量、但在引力上不与任何恒星或棕矮星发生束缚关系的星际天体$^{[1][2][3][4]}$。

流浪行星可能起源于其形成所在的行星系统，并在随后被抛射出去；它们也可能在行星系统之外独立形成。仅银河系中就可能存在数十亿到数万亿颗流浪行星，这一数量范围有望由即将投入使用的南希·格雷斯·罗曼空间望远镜进一步加以精确限定$^{[5][6]}$。流浪行星进入太阳系的概率极低，更不用说对地球生命构成直接威胁——在未来一千年内发生这种情况的概率约为一万亿分之一$^{[7]}$。

一些行星质量天体可能以类似恒星的方式形成，国际天文学联合会已提议将此类天体称为亚棕矮星（sub-brown dwarfs）$^{[8]}$。一个可能的例子是 Cha~110913$-$773444，它可能是在被抛射后成为一颗流浪行星，也可能是独立形成并成为一颗亚棕矮星$^{[9]}$。
\subsection{术语}
最早的两篇发现论文分别使用了“孤立的行星质量天体”（iPMO）$^{[10]}$和“自由漂浮行星”（FFP）$^{[11]}$这两个名称。大多数天文学论文使用其中之一$^{[12][13][14]}$。“流浪行星”这一术语更常用于引力微透镜研究，而这类研究也经常使用 FFP 这一术语$^{[15][16]}$。面向公众的新闻稿则可能采用不同的名称。例如，2021 年发现至少 70 颗 FFP 时，不同新闻稿分别使用了“流浪行星”$^{[17]}$、“无恒星行星”$^{[18]}$、“游荡行星”$^{[19]}$以及“自由漂浮行星”$^{[20]}$等称呼。
```
